\begin{frame}[t]{Success after reaching a checkpoint does not measure success from task start.\ev{\tagA}}
\vspace{0pt}
\begin{center}\begin{tikzpicture}[x=1cm,y=1cm,every node/.style={align=center,font=\fontsize{12.5}{15}\selectfont}]

\node[state,text width=3.4cm] (start) at (0,0) {Task start};
\node[candidate,text width=3.4cm] (cp) at (4.8,0) {Reach checkpoint\\probability 0.5};
\node[evidence,text width=3.4cm] (ok) at (9.6,0) {Succeed from there\\conditional probability 0.8};
\draw[flow](start)--(cp);\draw[flow](cp)--(ok);
\node[text=Teal,font=\fontsize{24}{28}\selectfont]at(4.8,-1.9){End-to-end success: $0.5\times0.8=0.4$};

\end{tikzpicture}\end{center}
\sources{Illustrative conditional-probability example; Go-Explore motivates returning to reached states.}
\qadetail{The values .5 and .8 are illustrative, not Go-Explore or factory measurements. Illustrative probabilities, not measured. A viable reached state contains information about the prefix; continuation success alone does not estimate how often it is reached or validate every earlier decision. Go-Explore explicitly separates return-to-state and exploration. A success rate at a selected checkpoint differs from a task-start success rate.}
\note{[0:45] A checkpoint changes the population we are studying. Suppose half of tasks reach a useful checkpoint, and eighty percent of continuations from that checkpoint succeed. End-to-end success is forty percent. Reporting eighty percent alone describes the selected reached-state population. Returning to a useful state can still be valuable; Go-Explore makes that strategy concrete. The evaluation must state whether it starts from tasks, from reached checkpoints, or from accepted checkpoints after a fidelity check. We now turn that distinction into acceptance obligations.}
\end{frame}
